% AA-15 | Algebra abstrakcyjna | Iteracja VIII
\chapter{Pierścienie nieprzemienne i algebry grupowe}
\label{aa:ch:noncommutative-rings}

\section{Dlaczego nie można zakładać przemienności?}
Dotychczas główne twierdzenia dotyczące ideałów, podzielności,
lokalizacji oraz pierścieni noetherowskich sformułowaliśmy dla
pierścieni przemiennych. Istnieją jednak naturalne pierścienie, w których
kolejność mnożenia ma znaczenie. Już dla przestrzeni wektorowej
$V$ nad ciałem $K$ pierścień $\operatorname{End}_K(V)$ z dodawaniem
odwzorowań i mnożeniem przez składanie jest zwykle nieprzemienny.
Nawet jeśli $K$ jest ciałem, macierze mogą nie komutować.

W tym rozdziale wszystkie pierścienie mają jedynkę $1_R\ne0_R$,
a homomorfizmy pierścieni zachowują jedynkę. Nie zakładamy
przemienności mnożenia, chyba że wyraźnie ją zaznaczamy.

\begin{definicja}[Pierścień niekoniecznie przemienny]
\label{aa:def:nc-ring}
Pierścień z jedynką to zbiór $R$ z działaniami $+$ i $\cdot$,
zerem $0_R$ i jedynką $1_R$ taki, że $(R,+,0_R)$ jest grupą
abelową, mnożenie jest łączne, $1_Ra=a1_R=a$ dla każdego
$a\in R$ i zachodzą obie rozdzielności:
$ a(b+c)=ab+ac$ oraz $(a+b)c=ac+bc$ dla wszystkich
$a,b,c\in R$. Pierścień jest przemienny, jeśli ponadto
$ab=ba$ dla wszystkich $a,b\in R$.
\end{definicja}

\begin{przyklad}[Macierze nie komutują]
\label{aa:prz:nc-matrix}
Niech $K$ będzie ciałem, a $E_{ij}$ oznacza macierz $2\times2$
mającą $1$ na pozycji $(i,j)$ i zera poza nią. Wówczas
$E_{12}E_{21}=E_{11}$, lecz $E_{21}E_{12}=E_{22}$.
Dlatego $M_2(K)$ nie jest pierścieniem przemiennym.
Jednocześnie obie macierze $E_{12}$ i $E_{21}$ są niezerowe,
a $E_{12}^2=E_{21}^2=0$. Przykład wyjaśnia, dlaczego
argumenty wykorzystujące skracanie czynników w dziedzinie
całkowitości nie przenoszą się bez dodatkowych założeń
na dowolne pierścienie macierzy.
\end{przyklad}

\begin{definicja}[Centrum i komutator elementów pierścienia]
\label{aa:def:nc-center}
Dla pierścienia $R$ definiujemy jego \emph{centrum}
$Z(R)=\{z\in R:\ za=az\text{ dla każdego }a\in R\}$.
Dla $a,b\in R$ ich \emph{komutatorem addytywnym}
nazywamy $[a,b]=ab-ba$. Nie należy utożsamiać tego
komutatora z komutatorem multiplikatywnym elementów grupy.
\end{definicja}

\begin{stwierdzenie}[Centrum jest podpierścieniem przemiennym]
\label{aa:st:nc-center-subring}
Niech $R$ będzie pierścieniem z jedynką. Wtedy $Z(R)$
jest przemiennym podpierścieniem zawierającym $1_R$.
\end{stwierdzenie}
\begin{proof}
Dla każdego $a\in R$ mamy $0_Ra=a0_R$ i $1_Ra=a1_R$.
Jeżeli $z,w\in Z(R)$ i $a\in R$, to
$(z-w)a=za-wa=az-aw=a(z-w)$ oraz
$(zw)a=z(wa)=z(aw)=(za)w=(az)w=a(zw)$.
Zatem $z-w,zw\in Z(R)$, co dowodzi zamkniętości
na działania i obecności jedynki. Dla $z,w\in Z(R)$
równość $zw=wz$ wynika z definicji centrum, ponieważ
$w$ także należy do $R$. Mnożenie w centrum jest
więc przemienne.
\end{proof}

\section{Ideały lewostronne, prawostronne i dwustronne}
\begin{definicja}[Trzy rodzaje ideałów]
\label{aa:def:nc-ideals}
Niech $R$ będzie pierścieniem. Addytywną podgrupę
$I\subseteq R$ nazywamy \emph{ideałem lewostronnym}, jeśli
$ri\in I$ dla każdego $r\in R$ i $i\in I$;
\emph{ideałem prawostronnym}, jeśli $ir\in I$;
\emph{ideałem dwustronnym}, jeśli zachodzą obie własności.
Dla $a\in R$ oznaczamy przez $Ra$, $aR$ i $RaR$
odpowiednio najmniejszy ideał lewostronny, prawostronny
i dwustronny zawierający $a$. Dokładniej,
$Ra=\{ra:r\in R\}$, $aR=\{ar:r\in R\}$ oraz
$RaR=\{\sum_{j=1}^n r_ja s_j:n\ge0,\ r_j,s_j\in R\}$,
gdzie dla $n=0$ suma ma wartość $0$.
Zbiory $Ra$ i $aR$ są addytywnymi podgrupami,
odpowiednio zamkniętymi na mnożenie z lewej lub prawej,
a dzięki jedynce zawierają $a$. Każdy odpowiedni ideał
zawierający $a$ zawiera także te zbiory.
Analogicznie $RaR$ jest grupą addytywną, zamkniętą
na mnożenie z lewej i prawej dzięki łączności mnożenia.
Zawiera $a=1_Ra1_R$, a każdy ideał dwustronny
zawierający $a$ zawiera wszystkie wyrazy $ras$ i ich
skończone sumy. To dowodzi minimalności $RaR$.
\end{definicja}

\begin{przyklad}[Ideał jednostronny nie musi być dwustronny]
\label{aa:prz:nc-one-sided}
Niech $R=M_2(K)$ i niech $I$ składa się z macierzy
mających drugą kolumnę równą zeru. Jest to ideał
lewostronny: dla $A\in R$ oraz $B\in I$ druga
kolumna $AB$ jest równa $A\cdot0=0$.
Nie jest jednak ideałem prawostronnym, ponieważ
$E_{11}\in I$, a $E_{11}E_{12}=E_{12}\notin I$.
Podobnie macierze o drugim wierszu równym zeru
tworzą ideał prawostronny: mnożenie z prawej
zachowuje zerowy drugi wiersz. Nie jest to ideał
lewostronny, gdyż $E_{11}$ do niego należy,
a $E_{21}E_{11}=E_{21}$ już nie.
Ograniczenie do ideałów dwustronnych nie jest zatem
jedynie zmianą nazwy: tylko one pozwalają utworzyć
pierścień ilorazowy.
\end{przyklad}

\begin{twierdzenie}[Kryterium istnienia pierścienia ilorazowego]
\label{aa:tw:nc-quotient}
Niech $R$ będzie pierścieniem z jedynką i $1_R\ne0_R$,
a $I$ jego addytywną podgrupą. Wzór
$(a+I)(b+I)=ab+I$ określa poprawnie mnożenie
klas w $R/I$ wtedy i tylko wtedy, gdy $I$ jest
ideałem dwustronnym. Wówczas $R/I$ ma łączne
mnożenie z jedynką $1_R+I$, a naturalny rzut
$R\to R/I$ zachowuje jedynkę. Przy przyjętej
w tym rozdziale konwencji $1\ne0$ iloraz jest
pierścieniem w tej samej klasie obiektów dokładnie wtedy,
gdy $I\ne R$. Dla $I=R$ otrzymujemy jednoelementowy
\emph{pierścień zerowy}, w którym $1=0$; jeśli
pierścienie zerowe są dopuszczone, także ten iloraz
jest pierścieniem.
\end{twierdzenie}
\begin{proof}
Załóżmy najpierw dwustronność $I$. Jeśli
$a'=a+i$ oraz $b'=b+j$ dla $i,j\in I$, to
$a'b'-ab=aj+ib+ij\in I$. Wynik mnożenia nie zależy
więc od reprezentantów. Łączność, rozdzielność oraz
jedynka $1_R+I$ przechodzą z $R$ na iloraz.
Przy tym $1_R+I=0_R+I$ jest równoważne
$1_R\in I$. Dla ideału dwustronnego zawierającego
jedynkę mamy $r=r1_R\in I$ dla każdego $r\in R$,
a więc $I=R$. To dowodzi także ostatniego
zastrzeżenia w tezie.

W przeciwną stronę załóżmy, że mnożenie klas
jest dobrze określone. Dla każdego $i\in I$ mamy
$i+I=0+I$. Jeśli $r\in R$, możemy pomnożyć tę
równość z lewej i z prawej strony przez $r+I$.
Dostajemy $ri+I=I$ i $ir+I=I$, a więc $ri,ir\in I$.
Zatem $I$ jest ideałem dwustronnym. Wzory na rzut
zachowują dodawanie, mnożenie i jedynkę z definicji.
\end{proof}

\begin{twierdzenie}[Pierwsze twierdzenie o izomorfizmie pierścieni]
\label{aa:tw:nc-first-isomorphism}
Niech $R,S$ będą pierścieniami z jedynką, a
$f:R\to S$ homomorfizmem zachowującym jedynkę.
Wówczas $\ker f$ jest ideałem dwustronnym i istnieje
kanoniczny izomorfizm $R/\ker f\cong f(R)$.
\end{twierdzenie}
\begin{proof}
Dla $x\in\ker f$ oraz $r\in R$ otrzymujemy
$f(rx)=f(r)f(x)=0$ i $f(xr)=f(x)f(r)=0$.
Jądro jest podgrupą addytywną, więc jest ideałem
dwustronnym. Określamy $\bar f(r+\ker f)=f(r)$.
Jeśli $r-r'\in\ker f$, to $f(r)=f(r')$, zatem
$\bar f$ jest poprawnie określonym homomorfizmem.
Obraz $f(R)$ jest podpierścieniem $S$: jest zamknięty
na różnice i iloczyny, a z warunku $f(1_R)=1_S$
wynika, że zawiera jedynkę $S$.
Odwzorowanie $\bar f$ jest surjektywne na ten podpierścień.
Ponadto $\bar f(r+\ker f)=0$ wtedy i tylko wtedy,
gdy $r\in\ker f$, czyli gdy $r+\ker f$ jest klasą
zerową. Stąd $\bar f$ jest iniektywne, a więc jest
izomorfizmem pierścieni z jedynką.
\end{proof}

\section{Pierścienie macierzy i ich ideały}
\begin{lemat}[Mnożenie jednostek macierzowych]
\label{aa:lem:nc-matrix-units}
Dla pierścienia $R$ z jedynką i jednostek macierzowych
$E_{ij}\in M_n(R)$, $n\ge1$, zachodzi
$E_{ij}E_{k\ell}=\delta_{jk}E_{i\ell}$, gdzie
$\delta_{jk}=1_R$, gdy $j=k$, oraz $\delta_{jk}=0_R$,
gdy $j\ne k$; macierz $E_{ij}$ ma $1_R$ w pozycji
$(i,j)$ i zera na pozostałych pozycjach.
\end{lemat}
\begin{proof}
Współczynnik na pozycji $(p,q)$ iloczynu wynosi
$\sum_{t=1}^n(E_{ij})_{pt}(E_{k\ell})_{tq}$.
Niezerowy składnik może wystąpić wyłącznie przy
$p=i$, $t=j=k$ i $q=\ell$; ma wtedy wartość $1_R$.
W pozostałych przypadkach suma jest zerowa.
\end{proof}

\begin{twierdzenie}[Dwustronne ideały pełnego pierścienia macierzy]
\label{aa:tw:nc-matrix-ideals}
Niech $R$ będzie pierścieniem z jedynką, a $n\ge1$.
Dla każdego ideału dwustronnego $J\triangleleft M_n(R)$
istnieje dokładnie jeden ideał dwustronny $I\triangleleft R$
taki, że $J=M_n(I)$, gdzie $M_n(I)$ oznacza zbiór
macierzy, których wszystkie współczynniki należą do $I$.
\end{twierdzenie}
\begin{proof}
Najpierw, jeśli $I$ jest ideałem dwustronnym $R$,
to $M_n(I)$ jest podgrupą addytywną. Dla
$A\in M_n(R)$ i $B\in M_n(I)$ współczynniki
$(AB)_{ij}=\sum_k a_{ik}b_{kj}$ oraz
$(BA)_{ij}=\sum_k b_{ik}a_{kj}$ należą do $I$.
Stąd $M_n(I)$ jest ideałem dwustronnym.

Niech teraz $J\triangleleft M_n(R)$. Niech $I$ będzie
zbiorem tych $a\in R$, dla których macierz $aE_{11}$
należy do $J$. Jest to podgrupa addytywna. Dla
$r\in R$ i $a\in I$ iloczyny
$(rE_{11})(aE_{11})=(ra)E_{11}$ oraz
$(aE_{11})(rE_{11})=(ar)E_{11}$ należą do $J$.
Zatem $I$ jest ideałem dwustronnym $R$.

Jeśli $B=(b_{ij})\in J$, to
$E_{1i}BE_{j1}=b_{ij}E_{11}\in J$: mnożenie przez
jednostki macierzowe wybiera współczynnik $b_{ij}$,
przy zachowaniu jego pozycji w iloczynie, również
gdy $R$ jest nieprzemienny. Zatem
$b_{ij}\in I$ dla wszystkich $i,j$. Dostaliśmy
$J\subseteq M_n(I)$. Odwrotnie, dla $a\in I$ mamy
$(aE_{11})\in J$, a zatem
$E_{i1}(aE_{11})E_{1j}=aE_{ij}\in J$. Każda
macierz z $M_n(I)$ jest sumą takich macierzy,
więc $M_n(I)\subseteq J$. Jednoznaczność $I$
wynika z warunku $aE_{11}\in J$.
\end{proof}

\begin{wniosek}[Prostota pierścienia macierzy nad ciałem]
\label{aa:wn:nc-matrix-simple}
Dla ciała $K$ i $n\ge1$ pierścień $M_n(K)$
nie posiada ideałów dwustronnych innych niż $(0)$
i cały pierścień, czyli jest \emph{pierścieniem prostym}.
Nie oznacza to, że jego ideały jednostronne są trywialne.
\end{wniosek}
\begin{proof}
Ciało $K$ ma dokładnie dwa ideały: $(0)$ i $K$.
Z Twierdzenia~\ref{aa:tw:nc-matrix-ideals} wynika,
że ideały dwustronne $M_n(K)$ mają postać
$M_n(0)$ lub $M_n(K)$. Dla $n\ge2$ przykład
z macierzami o jednej zerowej kolumnie daje
nietrywialny ideał lewostronny.
\end{proof}

\begin{twierdzenie}[Centrum pierścienia macierzy]
\label{aa:tw:nc-matrix-center}
Dla dowolnego pierścienia $R$ z jedynką i $n\ge1$
zachodzi $Z(M_n(R))=\{zI_n:z\in Z(R)\}$.
\end{twierdzenie}
\begin{proof}
Jeżeli $z\in Z(R)$, to dla każdej macierzy
$A=(a_{ij})$ zachodzi $(zI_n)A=(za_{ij})$ oraz
$A(zI_n)=(a_{ij}z)$, więc $zI_n$ należy do centrum.

Niech odwrotnie $B=(b_{ij})$ komutuje ze wszystkimi
macierzami. Komutując w szczególności z $E_{kk}$,
otrzymujemy $BE_{kk}=E_{kk}B$. Dla indeksów
$i\ne k$ i $j=k$ z tej równości wynika $b_{ik}=0$,
a dla $i=k$, $j\ne k$ wynika $b_{kj}=0$.
Ponieważ $k$ jest dowolne, macierz $B$ jest diagonalna.
Z równości $BE_{ij}=E_{ij}B$ dla $i\ne j$ wynika,
że wszystkie elementy diagonalne są jednakowe:
$b_{ii}=b_{jj}=z$. W przypadku $n=1$ jest to
prawdziwe bez dalszego argumentu. Wreszcie $B$
komutuje z macierzami $aE_{11}$ dla każdego $a\in R$,
co daje $za=az$. Tak więc $z\in Z(R)$ i $B=zI_n$.
\end{proof}


\begin{przyklad}[Ideały dwustronne a ideały jednostronne w macierzach]
\label{aa:prz:nc-matrix-one-sided-contrast}
Niech $K$ będzie ciałem i $R=M_2(K)$. Podzbiór
$I=\{\left(\begin{smallmatrix}a&0\\b&0\end{smallmatrix}\right):a,b\in K\}$
jest niezerowym właściwym ideałem lewostronnym. Istotnie,
$E_{11}\in I$, lecz $E_{22}\notin I$, a dla każdej macierzy
$A\in R$ i $B\in I$ druga kolumna $AB$ jest zerowa.
Nie jest to ideał dwustronny, gdyż $E_{11}E_{12}=E_{12}\notin I$.
Z drugiej strony Twierdzenie~\ref{aa:tw:nc-matrix-ideals} pokazuje,
że $R$ nie ma właściwych niezerowych ideałów dwustronnych.
Nie można więc zastąpić słowa „dwustronny” słowem „lewostronny”
w twierdzeniu o prostocie pierścienia macierzy.
\end{przyklad}

\section{Pierścień przeciwny i moduły prawostronne}
\begin{definicja}[Pierścień przeciwny]
\label{aa:def:nc-opposite}
Dla pierścienia $R$ definiujemy pierścień przeciwny
$R^{\mathrm{op}}$ o tej samej grupie addytywnej i
mnożeniu $a\star b=ba$, gdzie po prawej mnożymy w $R$.
Jedynka pozostaje ta sama. Łączność jest zachowana,
ponieważ $(a\star b)\star c=c(ba)=(cb)a=
a\star(b\star c)$; obie rozdzielności wynikają
z rozdzielności w $R$ po odwróceniu kolejności
czynników.
\end{definicja}
\begin{stwierdzenie}[Zmiana strony działania]
\label{aa:st:nc-opposite-module}
Prawy $R$-moduł $M$ jest równoważnie lewym
$R^{\mathrm{op}}$-modułem poprzez działanie
$a\star m:=ma$. Ta zgodność działa także dla
homomorfizmów modułów.
\end{stwierdzenie}
\begin{proof}
Prawy moduł spełnia $(ma)b=m(ab)$.
W $R^{\mathrm{op}}$ mamy $a\star b=ba$,
a po stronie modułu
$(a\star b)\star m=m(ba)$.
Z drugiej strony
$a\star(b\star m)=(mb)a=m(ba)$, więc
łączność działania jest zachowana. Rozdzielności
oraz działanie jedynki przechodzą bezpośrednio
z aksjomatów prawego modułu. Homomorfizm, który
zachowuje prawostronne mnożenie przez skalary,
zachowuje z definicji odpowiadające mu lewostronne.
\end{proof}

\section{Pierścienie grupowe i odwzorowania uniwersalne}
\begin{definicja}[Pierścień grupowy]
\label{aa:def:nc-group-ring}
Niech $R$ będzie przemiennym pierścieniem z jedynką,
a $G$ dowolną grupą. \emph{Pierścień grupowy} $R[G]$
to zbiór wszystkich formalnych skończonych sum
$\sum_{g\in G}r_g[g]$, gdzie $r_g\in R$, a niemal
wszystkie współczynniki są zerowe. Dodawanie jest
współrzędne, mnożenie określa się przez
\[
\left(\sum_g r_g[g]\right)\left(\sum_h s_h[h]\right)
 =\sum_{k\in G}\left(\sum_{gh=k}r_gs_h\right)[k].
\]
Suma wewnętrzna jest skończona, ponieważ obie rodziny
współczynników mają skończony nośnik. Jedynką jest
$1_R[1_G]$. W przypadku $R=K$, gdzie $K$ jest
ciałem, pierścień $K[G]$ nazywa się również
\emph{algebrą grupową} grupy $G$ nad $K$.
\end{definicja}
\begin{stwierdzenie}[Aksjomaty pierścienia grupowego]
\label{aa:st:nc-group-ring-ring}
Dla przemiennego pierścienia $R$ z jedynką i dowolnej
grupy $G$ powyższe działania czynią $R[G]$ pierścieniem.
Ponadto $[g][h]=[gh]$ oraz $[g]^{-1}=[g^{-1}]$.
Jeżeli $R\ne0$, to $R[G]$ jest przemienny wtedy
i tylko wtedy, gdy $G$ jest abelowa.
\end{stwierdzenie}
\begin{proof}
Skończone nośniki gwarantują, że każdy współczynnik
iloczynu jest sumą skończoną. Łączność wynika z
obliczenia współczynnika przy $[k]$ dla trzech sum:
zarówno $(xy)z$, jak i $x(yz)$ mają współczynnik
$\sum_{abc=k}r_as_bt_c$. Jest to suma skończona,
bo indeksy $a,b,c$ należą do skończonych nośników
trzech czynników. Zmiana nawiasowania jest
więc poprawna dzięki łączności mnożenia w $R$
i w $G$. Rozdzielności sprawdzamy współczynnikowo,
a element $1_R[1_G]$ zachowuje wszystkie współczynniki.
To dowodzi aksjomatów pierścienia; tożsamości
$[g][h]=[gh]$ i $[g][g^{-1}]=[1_G]$ są bezpośrednie.
Jeśli $G$ jest abelowa i $R$ przemienny, wzór na
mnożenie pokazuje, że zamiana kolejności czynników
nie zmienia iloczynu. Odwrotnie, jeśli $R[G]$
jest przemienny, to $[gh]=[hg]$. Ponieważ
$R\ne0$, formalne symbole $[g]$ są różnymi
elementami wolnego modułu dla różnych $g$.
Zatem $gh=hg$, czyli $G$ jest abelowa.
\end{proof}

\begin{twierdzenie}[Własność uniwersalna pierścienia grupowego]
\label{aa:tw:nc-group-ring-universal}
Niech $R$ będzie przemiennym pierścieniem z jedynką,
a $G$ grupą. Niech $S$ będzie pierścieniem z jedynką,
$f:R\to S$ homomorfizmem pierścieni takim, że
$f(R)\subseteq Z(S)$, i niech $u:G\to S^\times$
będzie homomorfizmem grup do grupy jednostek $S$.
Wtedy istnieje dokładnie jeden homomorfizm pierścieni
$F:R[G]\to S$ spełniający $F(r[1_G])=f(r)$
oraz $F([g])=u(g)$ dla wszystkich $r\in R$, $g\in G$.
Jest nim odwzorowanie
$F(\sum_g r_g[g])=\sum_g f(r_g)u(g)$.
\end{twierdzenie}
\begin{proof}
Ponieważ zapisy w $R[G]$ mają jednoznaczne współczynniki,
podany wzór określa odwzorowanie. Zachowuje dodawanie,
bo sumowanie jest współrzędne. Dla dwóch składników
$r[g]$, $s[h]$ mamy
\[
F((r[g])(s[h]))=f(rs)u(gh)
=f(r)f(s)u(g)u(h).
\]
Warunek $f(R)\subseteq Z(S)$ pozwala przestawić
$f(s)$ z $u(g)$, więc prawa strona jest równa
$f(r)u(g)f(s)u(h)=F(r[g])F(s[h])$.
Rozdzielność i skończoność sum rozszerzają to na
dowolne elementy. Zachowanie jedynki daje
$F(1_R[1_G])=f(1_R)u(1_G)=1_S$.
Jeśli $H$ jest innym homomorfizmem o tych wartościach,
to z równości $r[g]=(r[1_G])[g]$ i addytywności
wynika $H(\sum r_g[g])=\sum f(r_g)u(g)$.
Zatem $H=F$.
\end{proof}

\begin{przyklad}[Grupa cykliczna i algebra wielomianów]
\label{aa:prz:nc-cyclic-ring}
Niech $n\ge1$ będzie liczbą całkowitą, $K$ ciałem,
a $C_n=\langle g\rangle$ grupą cykliczną rzędu
$n$ (czyli $g^n=1$ i nie istnieje mniejszy dodatni
wykładnik z tą własnością). Wówczas otrzymujemy izomorfizm
$K[C_n]\cong K[t]/(t^n-1)$. Homomorfizm
$K[t]\to K[C_n]$ dany przez $t\mapsto[g]$
jest surjektywny, gdyż każda potęga $[g]^j$
pochodzi od $t^j$. Ponieważ $[g]^n=[1]$,
wielomian $t^n-1$ należy do jądra. Indukowany
homomorfizm $K[t]/(t^n-1)\to K[C_n]$ jest więc
surjektywny. Po lewej każdy element ma jednoznacznego
reprezentanta stopnia mniejszego niż $n$, zatem
wymiar nad $K$ wynosi $n$. Po prawej bazą są
$[1],[g],\ldots,[g^{n-1}]$, więc wymiar także wynosi
$n$. Surjektywne odwzorowanie liniowe między
przestrzeniami tego samego skończonego wymiaru
jest izomorfizmem.
\end{przyklad}

\begin{przyklad}[Grupa nieabelowa daje pierścień nieprzemienny]
\label{aa:prz:nc-s3-ring}
Niech $K$ będzie dowolnym ciałem i $G=S_3$.
Elementy $[(12)]$ i $[(23)]$ są jednostkami pierścienia
$K[S_3]$, ale nie komutują, ponieważ
$(12)(23)\ne(23)(12)$ w $S_3$.
Jednocześnie $K[S_3]$ ma wymiar sześć nad $K$:
bazą jest sześć formalnych symboli $[g]$.
To prosty przykład sześciowymiarowej algebry
nieprzemiennej powstałej bez udziału macierzy.
\end{przyklad}


\begin{przyklad}[Nilpotenty w pierścieniu grupowym]
\label{aa:prz:nc-group-ring-nilpotent}
Niech $p$ będzie liczbą pierwszą, $K=\mathbb F_p$,
a $G=C_p=\langle g\rangle$. W pierścieniu $K[G]$
element $u=[g]-[1_G]$ jest niezerowy, ponieważ $g\ne1_G$
i formalne symbole $[g]$, $[1_G]$ są liniowo niezależne nad $K$.
Ponieważ $[g]$ i $[1_G]$ komutują, dwumian Newtona daje
$u^p=[g]^p-[1_G]^p=[g^p]-[1_G]=0$:
wszystkie pośrednie współczynniki dwumianowe są podzielne
przez $p$. Zatem nawet pierścień grupowy nad ciałem
może zawierać niezerowe elementy nilpotentne. Założenie
$\operatorname{char}K=p$ jest istotne: dla $K=\mathbb Q$
i $G=C_p$ w rzeczywistości $u^p\ne0$.
Istotnie, we wzorze dwumianowym współczynnik przy
$[g]$ wynosi $(-1)^{p-1}p$, ponieważ jedynymi
potęgami $g$ równymi $g$ wśród $g^0,\ldots,g^p$
są $g^1$. Współczynnik ten jest niezerowy w
$\mathbb Q$, więc $u^p\ne0$.
\end{przyklad}

\section{Wolne pierścienie nieprzemienne}
\begin{definicja}[Pierścień wolnych wielomianów]
\label{aa:def:nc-free-associative}
Niech $K$ będzie ciałem, a $X$ zbiorem. Pierścień
$K\langle X\rangle$ składa się ze skończonych
formalnych kombinacji $\sum_w a_ww$ słów $w$ w
alfabecie $X$, włącznie ze słowem pustym, będącym
jedynką. Iloczyn słów to ich konkatenacja,
a iloczyn sum określamy przez rozdzielność.
Skalary z $K$ komutują ze wszystkimi słowami.
\end{definicja}

\begin{twierdzenie}[Własność uniwersalna wolnego pierścienia]
\label{aa:tw:nc-free-universal}
Niech $K$ będzie ciałem, $X$ zbiorem, a $S$ pierścieniem
z jedynką wyposażonym w homomorfizm $f:K\to Z(S)$.
Dla każdej funkcji $u:X\to S$ istnieje dokładnie jeden
homomorfizm pierścieni $F:K\langle X\rangle\to S$
taki, że $F(a)=f(a)$ dla $a\in K$ oraz
$F(x)=u(x)$ dla $x\in X$.
\end{twierdzenie}
\begin{proof}
Wartość $F$ na słowie $x_1\cdots x_m$ jest wymuszona
przez zachowanie iloczynu i wynosi
$u(x_1)\cdots u(x_m)$; słowo puste musi przejść na
$1_S$. Rozszerzamy ten przepis $K$-liniowo:
$F(\sum_w a_ww)=\sum_w f(a_w)F(w)$.
Skończoność nośnika zapewnia poprawność definicji.
Dla słów $v,w$ zachodzi $F(vw)=F(v)F(w)$ z
łączności w $S$. Dla kombinacji liniowych własność
wynika z rozdzielności; warunek centralności $f(K)$
gwarantuje, że współczynniki można przesuwać
przez obrazy słów. Jednoznaczność wynika z tego,
że elementy $X$ wraz ze skalarami generują cały
pierścień $K\langle X\rangle$.
\end{proof}

\begin{przyklad}[Dlaczego wielomiany nieprzemienne są potrzebne?]
\label{aa:prz:nc-free-example}
Dla $X=\{x,y\}$ słowa $xy$ i $yx$ są różnymi
wektorami bazowymi $K\langle x,y\rangle$; dlatego
$xy-yx\ne0$. W pierścieniu przemiennym $K[x,y]$
ten sam formalny komutator znika. Kanoniczny
homomorfizm $K\langle x,y\rangle\to K[x,y]$
utożsamia dokładnie relacje wymuszające przemienność
$x$ i $y$: jego jądrem jest ideał dwustronny
$J=(xy-yx)$ generowany przez $xy-yx$.
Wykażemy obie inkluzje przez skonstruowanie
izomorfizmu. Odwzorowanie $\pi:K\langle x,y\rangle
\to K[x,y]$, przypisujące słowu przemienny
jednomian o tych samych wykładnikach $x,y$,
jest surjektywnym homomorfizmem, a $\pi(xy-yx)=0$.
Stąd $J\subseteq\ker\pi$ i dostajemy surjekcję
$\bar\pi:Q\to K[x,y]$, gdzie
$Q=K\langle x,y\rangle/J$.
W pierścieniu $Q$ klasy $\bar x,\bar y$ komutują.
Określmy $\psi:K[x,y]\to Q$ wzorem
$\psi(\sum_{a,b}c_{ab}x^ay^b)=
\sum_{a,b}c_{ab}\bar x^a\bar y^b$.
Wzór jest poprawnie określony dzięki jednoznacznemu
zapisowi wielomianu oraz wyznacza homomorfizm,
ponieważ $\bar x\bar y=\bar y\bar x$.
Ponadto $\bar\pi\circ\psi$ jest identycznością
na $K[x,y]$. Z kolei $\psi\circ\bar\pi$ jest
identycznością na $Q$: obie strony zgadzają się
na klasach $x,y$ i skalarach z $K$, które generują
cały $Q$ jako pierścień.
Można to sprawdzić także bezpośrednio na słowach:
w $Q$ każda zamiana sąsiednich liter $yx$ na $xy$
nie zmienia klasy słowa, ponieważ dla dowolnych słów
$v,w$ różnica $v(yx-xy)w$ należy do $J$.
Powtarzanie takich zamian ściśle zmniejsza liczbę par,
w których litera $y$ poprzedza literę $x$.
Po skończenie wielu krokach każde słowo o $a$ literach
$x$ i $b$ literach $y$ przechodzi więc w $x^ay^b$.
Zatem $\bar\pi$ jest izomorfizmem i $\ker\pi=J$,
co dowodzi tezy.
To pokazuje formalną różnicę między pierścieniem
przemiennym a nieprzemiennym.
\end{przyklad}
